\documentclass[12pt,a4paper]{article}

% ══════════════════════════════════════════════════════════════
% Preâmbulo compartilhado pelos documentos da P2
% Reaproveita o estilo de docs/documento-unificado.tex
% ══════════════════════════════════════════════════════════════

\usepackage[utf8]{inputenc}
\usepackage[T1]{fontenc}
\usepackage[brazil]{babel}

\usepackage[top=2.5cm, bottom=2.5cm, left=3cm, right=2cm]{geometry}
\usepackage{setspace}
\onehalfspacing
\setlength{\parindent}{0pt}
\setlength{\parskip}{6pt}

\usepackage{lmodern}
\usepackage{microtype}
\usepackage{amssymb}
\usepackage{amsmath}
\usepackage{pifont}
\newcommand{\cmark}{\ding{51}}
\newcommand{\xmark}{\ding{55}}

\usepackage[table,dvipsnames]{xcolor}
\definecolor{accent}{HTML}{2C7A4B}
\definecolor{accentdim}{HTML}{1E5934}
\definecolor{accentlight}{HTML}{E8F4ED}
\definecolor{tabheader}{HTML}{2C3E50}
\definecolor{tabrow}{HTML}{F5F6FA}
\definecolor{linkblue}{HTML}{2980B9}
\definecolor{codebg}{HTML}{F5F5F5}
\definecolor{codeborder}{HTML}{D0D0D0}
\definecolor{ndvilow}{HTML}{8B4513}
\definecolor{ndvimid}{HTML}{DAA520}
\definecolor{ndvihigh}{HTML}{006400}
\definecolor{warnyellow}{HTML}{F39C12}
\definecolor{successgreen}{HTML}{27AE60}
\definecolor{errorred}{HTML}{E74C3C}

\usepackage{booktabs}
\usepackage{tabularx}
\usepackage{multirow}
\usepackage{makecell}
\usepackage{float}
\usepackage{graphicx}
\usepackage{caption}
\usepackage{array}

\usepackage{enumitem}
\setlist[itemize]{itemsep=2pt, parsep=0pt, topsep=4pt, leftmargin=*}
\setlist[enumerate]{itemsep=2pt, parsep=0pt, topsep=4pt, leftmargin=*}

\usepackage[colorlinks=true, linkcolor=linkblue, urlcolor=linkblue, citecolor=linkblue]{hyperref}

\usepackage{fancyhdr}
\pagestyle{fancy}
\fancyhf{}
\fancyhead[L]{\small\textcolor{gray}{EasyGis}}
\fancyhead[R]{\small\textcolor{gray}{P2 — Gerência e Manutenção de Software · Prof.\ Mauro · 2026.1}}
\fancyfoot[C]{\thepage}
\renewcommand{\headrulewidth}{0.4pt}

\usepackage{titlesec}
\titleformat{\section}
  {\Large\bfseries\color{accentdim}}
  {\thesection.}{0.5em}{}
\titleformat{\subsection}
  {\large\bfseries\color{accent}}
  {\thesubsection}{0.5em}{}
\titleformat{\subsubsection}
  {\normalsize\bfseries}
  {\thesubsubsection}{0.5em}{}

% Listings (código)
\usepackage{listings}
\lstdefinestyle{shellstyle}{
  basicstyle=\ttfamily\small,
  backgroundcolor=\color{codebg},
  frame=single,
  rulecolor=\color{codeborder},
  framesep=4pt,
  xleftmargin=8pt,
  xrightmargin=4pt,
  breaklines=true,
  showstringspaces=false,
  keepspaces=true,
  literate={─}{{-}}1 {├}{{|}}1 {│}{{|}}1 {└}{{`}}1
}
\lstdefinestyle{pythonstyle}{
  basicstyle=\ttfamily\footnotesize,
  language=Python,
  backgroundcolor=\color{codebg},
  frame=single,
  rulecolor=\color{codeborder},
  framesep=4pt,
  xleftmargin=8pt,
  xrightmargin=4pt,
  breaklines=true,
  keywordstyle=\color{accent}\bfseries,
  stringstyle=\color{ndvimid},
  commentstyle=\color{gray}\itshape,
  showstringspaces=false,
  keepspaces=true
}
\lstset{style=shellstyle}

% Caixas para código inline
\newcommand{\code}[1]{\texttt{\colorbox{codebg}{\strut#1}}}

% Caixas de destaque
\usepackage{tcolorbox}
\tcbuselibrary{skins, breakable}

\newtcolorbox{infobox}[1][]{%
  colback=accentlight, colframe=accent,
  boxrule=0.6pt, arc=2pt,
  left=8pt, right=8pt, top=4pt, bottom=4pt,
  fonttitle=\bfseries, breakable, #1
}

\newtcolorbox{warnbox}[1][]{%
  colback=yellow!8, colframe=warnyellow,
  boxrule=0.6pt, arc=2pt,
  left=8pt, right=8pt, top=4pt, bottom=4pt,
  breakable, #1
}

\newtcolorbox{successbox}[1][]{%
  colback=green!4, colframe=successgreen,
  boxrule=0.6pt, arc=2pt,
  left=8pt, right=8pt, top=4pt, bottom=4pt,
  breakable, #1
}


\title{\textbf{Perguntas Prováveis · Apresentação P2}\\[6pt]
       \large EasyGis — Banco de respostas preparadas}
\author{Engenharia de Software · 2026.1}
\date{Maio · 2026}

\begin{document}
\maketitle

\section*{Como usar este documento}

As perguntas estão agrupadas em \textbf{8 blocos temáticos}. Cada
pergunta tem:

\begin{itemize}[noitemsep]
  \item \textbf{Q:} a pergunta literal que o professor pode fazer.
  \item \textbf{A:} resposta preparada, curta e direta.
  \item \textbf{Quem responde:} o integrante com domínio do assunto.
\end{itemize}

\textbf{Estratégia geral:}
\begin{itemize}
  \item Se a pergunta cair no seu slide, responda você. Se for sobre área
        de outro, passe a bola: \textquotedblleft\,O X tem mais contexto sobre isso.\,\textquotedblright
  \item \textbf{Nunca chute.} Se não souber, diga \textquotedblleft\,Não
        documentamos isso ainda, mas está em [doc X]\,\textquotedblright\
        ou \textquotedblleft\,Posso voltar com a resposta após a
        sessão.\,\textquotedblright
  \item \textbf{Cite o documento} sempre que possível
        (\texttt{00-ENTREGA-P2.pdf}, \texttt{04-diagramas-arquitetura.md}, etc.).
\end{itemize}

\newpage

\section*{Bloco 1 — Arquitetura}

\subsection*{Q1.1 — Por que dois processos separados e não monólito?}
\textbf{A:} Cada lado fala a linguagem ideal: Python tem o ecossistema
geoespacial (rasterio, NumPy, GDAL); JavaScript renderiza mapa e 3D em WebGL.
Acoplar tudo em um único processo nos forçaria a escolher um lado e perder o
outro. REST + JSON é o contrato mais simples e debugável entre os dois mundos.
\\[2pt] \textit{Responde: Vinicius Henrique.}

\subsection*{Q1.2 — Por que não tem banco de dados?}
\textbf{A:} Decisão consciente, registrada em \textbf{DA-02}. O escopo do MVP
é single-user em localhost, processando arquivos \texttt{.SAFE} e KMLs já no
filesystem. Adicionar banco aumentaria complexidade sem benefício imediato.
Está no roadmap (próximo passo 02: SQLite ou PostGIS) quando virar
multi-usuário.
\\[2pt] \textit{Responde: Vinicius Henrique.}

\subsection*{Q1.3 — Por que imagem em base64 dentro do JSON, e não um endpoint que retorna binário?}
\textbf{A:} Por simplicidade do contrato: uma resposta única
(\texttt{IndexResult}) traz a imagem, as estatísticas e o histograma juntos.
O frontend processa em um só lugar. O custo do base64 é ~33\% de overhead,
aceitável para imagens de talhão (~100\,KB). Se virasse gargalo, trocaríamos
para endpoint binário separado.
\\[2pt] \textit{Responde: João Leonardi.}

\subsection*{Q1.4 — O sistema é escalável?}
\textbf{A:} Hoje não — é single-user em localhost. Mas a separação
frontend/backend permite escalar horizontalmente o backend FastAPI atrás de
um load balancer quando precisar. O processador é sem estado, então
escala trivialmente.
\\[2pt] \textit{Responde: Vinicius Henrique.}

\section*{Bloco 2 — Stack tecnológica}

\subsection*{Q2.1 — Por que FastAPI e não Django ou Flask?}
\textbf{A:} FastAPI dá validação automática via Pydantic e gera OpenAPI em
\texttt{/docs} sem código extra. Django é overkill (ORM, admin, autenticação
nativos) para o que é essencialmente um endpoint de cálculo. Flask exigiria
plugins (\texttt{flask-pydantic}, \texttt{flasgger}) para o mesmo nível.
\\[2pt] \textit{Responde: Luis Gustavo.}

\subsection*{Q2.2 — Por que Leaflet e não Mapbox / Cesium?}
\textbf{A:} Leaflet é maduro, leve e \textbf{não exige token pago}. Mapbox
GL cobra após cota. Cesium é overkill para terreno 2D — usamos Three.js
diretamente para o 3D, com mais controle.
\\[2pt] \textit{Responde: Luis Gustavo.}

\subsection*{Q2.3 — Por que Next.js e não React puro?}
\textbf{A:} O App Router do Next.js permite ler KMLs server-side via API
routes (\texttt{/api/fields}), sem expor o filesystem ao cliente. Em React
puro precisaríamos de um servidor separado só pra isso.
\\[2pt] \textit{Responde: Luis Gustavo.}

\subsection*{Q2.4 — Por que \texttt{uv} e não pip + requirements.txt?}
\textbf{A:} \texttt{uv} é cerca de 10\,× mais rápido em instalação, gera lock
file determinístico (\texttt{uv.lock}) e simplifica setup com um único comando
\texttt{uv sync}. Trade-off: dependência extra a instalar — mitigado pela
documentação no README.
\\[2pt] \textit{Responde: Luis Gustavo.}

\section*{Bloco 3 — Backlog e priorização}

\subsection*{Q3.1 — Como vocês priorizaram o que entrar e o que ficar de fora?}
\textbf{A:} Aplicamos MoSCoW. Tudo que era \textbf{Must Have} (núcleo de
processamento — RF-01 a 11) entrou primeiro. Os 5 \textbf{Won't Have}
(autenticação, persistência em banco, exportação, multi-usuário, integração
Copernicus) foram excluídos do MVP \textbf{conscientemente}: o objetivo era
validar o núcleo geoespacial em um semestre, não construir um SaaS completo.
\\[2pt] \textit{Responde: José Melquíades.}

\subsection*{Q3.2 — Por que 30/35 e não 35/35?}
\textbf{A:} Das 35 histórias: 28 entregues completas + 2 parciais
(download manual de imagem) + 5 conscientemente fora (Won't Have). Os 5
fora não falharam — \textbf{foram planejados como fora desde o início}.
85,7\% é a entrega real do que prometemos.
\\[2pt] \textit{Responde: José Melquíades.}

\subsection*{Q3.3 — Como vocês garantem que o backlog está rastreado para os testes?}
\textbf{A:} Matriz RF\,×\,US em \texttt{03-especificacao-requisitos-SRS.md
§\,6}. Cada caso de teste em \texttt{tests/} cita o RF coberto no docstring.
RF-04, 05, 06, 10, 11, 20 têm cobertura direta.
\\[2pt] \textit{Responde: Lauan Matheus.}

\section*{Bloco 4 — Testes}

\subsection*{Q4.1 — A cobertura global de 24\% não é baixa demais?}
\textbf{A:} À primeira vista sim, mas é \textbf{honesto e explicado}. Esse
número é puxado para baixo pelos endpoints HTTP, que exigem um produto
Sentinel-2 \texttt{.SAFE} real (1\,GB) para serem exercitados. Onde a lógica
de cálculo realmente vive — no \texttt{sentinel\_processor.py} — a cobertura
é \textbf{55\%}. Declaramos o gap como \textbf{DT-03} no documento de entrega.
\\[2pt] \textit{Responde: Lauan Matheus.}

\subsection*{Q4.2 — O que vocês fariam para subir a cobertura?}
\textbf{A:} Três coisas: \textbf{(1)} mockar o rasterio nos testes de
endpoint para não depender de \texttt{.SAFE}; \textbf{(2)} criar fixtures de
produto Sentinel-2 mínimo (uma cena recortada de poucos KB);
\textbf{(3)} CI com GitHub Actions rodando a suíte em cada PR.
\\[2pt] \textit{Responde: Lauan Matheus.}

\subsection*{Q4.3 — O que são as \textquotedblleft fixtures sintéticas em NumPy\textquotedblright?}
\textbf{A:} Em vez de carregar um \texttt{.SAFE} real, geramos arrays NumPy
com o formato esperado (\texttt{shape=(H,W)}, \texttt{dtype=float32}) e
valores conhecidos. Permite testar NDVI, EVI, máscaras e estatísticas em
milissegundos, sem I/O de disco.
\\[2pt] \textit{Responde: Lauan Matheus.}

\subsection*{Q4.4 — Tem teste de UI / E2E automatizado?}
\textbf{A:} Hoje não. O \textbf{Nível 3} (manual) cobre o fluxo end-to-end
durante a demo. Adicionar Playwright ou Cypress está no backlog futuro — foi
deixado fora porque o custo de manter testes de UI brittle não compensava no
tempo do semestre.
\\[2pt] \textit{Responde: Lauan Matheus.}

\section*{Bloco 5 — Algoritmos e processamento}

\subsection*{Q5.1 — Por que NDVI e não outro índice como padrão?}
\textbf{A:} NDVI é o índice mais estabelecido e validado em sensoriamento
remoto agrícola — comparável a literatura, software comercial e produtos
oficiais. EVI e SAVI estão implementados também (RF-06), mas NDVI é o
\textit{default} por familiaridade do usuário.
\\[2pt] \textit{Responde: João Leonardi.}

\subsection*{Q5.2 — Como vocês tratam pixels de nuvem ou água?}
\textbf{A:} Hoje a máscara binária do polígono recorta apenas o talhão.
Pixels de nuvem dentro do polígono entram no cálculo — limitação consciente.
A próxima evolução é usar a banda QA60 (\textit{cloud mask}) do Sentinel-2
para excluir esses pixels.
\\[2pt] \textit{Responde: João Leonardi.}

\subsection*{Q5.3 — Por que a classificação K-Means e não Random Forest?}
\textbf{A:} K-Means é \textbf{não-supervisionado} — funciona sem rótulos
de treinamento, que não temos para Picos-PI. Random Forest exigiria um
dataset rotulado por agrônomo, fora do escopo do MVP. Está no roadmap
(próximo passo 03) quando tivermos dados rotulados.
\\[2pt] \textit{Responde: João Leonardi.}

\subsection*{Q5.4 — O cálculo de área em hectares está correto?}
\textbf{A:} Sim. Usamos a fórmula do \textbf{laço} (\textit{shoelace}) sobre
as coordenadas reprojetadas em UTM — sistema métrico nativo. Sem
reprojeção, o cálculo em graus daria área inflada. Verificável contra área
no Google Earth.
\\[2pt] \textit{Responde: João Leonardi.}

\section*{Bloco 6 — Demo / produto}

\subsection*{Q6.1 — Que tipo de imagem o sistema aceita?}
\textbf{A:} Produtos Sentinel-2 Level-2A no formato \texttt{.SAFE} (estrutura
padrão da ESA). O usuário hoje baixa do Copernicus Hub manualmente; integração
direta com a API está no roadmap (próximo passo 01).
\\[2pt] \textit{Responde: João V. Castello.}

\subsection*{Q6.2 — Quanto tempo leva pra processar uma imagem?}
\textbf{A:} \textbf{5 a 30 segundos} dependendo do tamanho do talhão e do
índice. NDVI é o mais rápido (duas bandas); EVI usa três e é ligeiramente
mais lento. A suavização gaussiana custa cerca de 1\,s extra em talhões grandes.
\\[2pt] \textit{Responde: João V. Castello.}

\subsection*{Q6.3 — Posso desenhar talhões ou só importar KML?}
\textbf{A:} Os dois. A versão atual aceita KML do Google Earth e
também tem ferramenta de desenho de polígono no mapa Leaflet. O desenho
não persiste entre sessões (limitação 1 do slide 13) — está no roadmap.
\\[2pt] \textit{Responde: João V. Castello.}

\subsection*{Q6.4 — E se a imagem tiver erro? Quebra a interface?}
\textbf{A:} Não. O backend faz validação Pydantic da entrada e retorna
\textbf{422 / 404 / 500} estruturado. O frontend exibe um toast com a
mensagem de erro. Cobertura desse caminho está no nível de integração
(FastAPI TestClient).
\\[2pt] \textit{Responde: João Leonardi.}

\subsection*{Q6.5 — A interface está acessível?}
\textbf{A:} Componentes shadcn/ui (Radix por baixo) garantem WAI-ARIA básico
— foco visível, navegação por teclado, contraste mínimo. Não rodamos
auditoria Lighthouse formal — é uma melhoria de roadmap.
\\[2pt] \textit{Responde: João V. Castello.}

\section*{Bloco 7 — Processo e equipe}

\subsection*{Q7.1 — Como vocês dividiram o trabalho?}
\textbf{A:} Por afinidade técnica e área de especialização, registrada no
\textbf{slide 14}: backend (João Leonardi), frontend (João Castello), dados
(José), testes (Lauan), pitch (Lucas), docs (Luis Gustavo), manual (Sammuel),
arquitetura (Vinicius). Cada um teve ownership claro de sua área e código
revisado por par em PR.
\\[2pt] \textit{Responde: João Leonardi.}

\subsection*{Q7.2 — Usaram alguma metodologia ágil?}
\textbf{A:} Sprints curtos com ciclo semanal informal — reuniões de
alinhamento e merge de PRs. Sem Scrum formal (sem PO/SM definidos), mas com
backlog priorizado e iteração visível no \texttt{git log}.
\\[2pt] \textit{Responde: João Leonardi.}

\subsection*{Q7.3 — Quanto tempo a equipe gastou no projeto?}
\textbf{A:} Aproximadamente um semestre letivo. As 28 histórias entregues
estão distribuídas em ~50 commits. Histórico completo em \texttt{git log}.
\\[2pt] \textit{Responde: João Leonardi.}

\subsection*{Q7.4 — Tem code review?}
\textbf{A:} Sim. Todas as features foram mergeadas via PR no GitHub
(\texttt{luisg0c/EasyGis}). Cada PR tinha pelo menos um aprovador além do
autor. Sem CI bloqueando merge — foi limitação consciente
(próximo passo 07).
\\[2pt] \textit{Responde: João Leonardi.}

\section*{Bloco 8 — Documentação e decisões de design}

\subsection*{Q8.1 — O que é um ADR?}
\textbf{A:} Architecture Decision Record. É um registro curto que captura
\textbf{contexto, decisão e consequências} de uma escolha técnica
importante. Documentamos 11 deles, de DA-01 (estrutura geral) a DA-11
(\texttt{uv} em vez de pip). Localizados em
\texttt{04-diagramas-arquitetura.md §\,9} e \texttt{00-ENTREGA-P2.pdf §\,3}.
\\[2pt] \textit{Responde: Luis Gustavo.}

\subsection*{Q8.2 — Qual decisão técnica mais difícil vocês tomaram?}
\textbf{A:} A da \textbf{DA-09} — suavização gaussiana preservando NaN.
Aplicar gaussiano direto sangrava valores para fora do polígono criando
halo. A correção pela máscara filtrada é não-óbvia e é o que distingue um
resultado correto de um plausível-mas-errado.
\\[2pt] \textit{Responde: Luis Gustavo ou João Leonardi.}

\subsection*{Q8.3 — Vocês consideraram alternativas para a stack atual?}
\textbf{A:} Sim. Considerei Django + DRF (rejeitada por overkill); Mapbox
(rejeitada por custo); Cesium para 3D (rejeitada por complexidade);
pip puro (rejeitada por velocidade). Cada rejeição está documentada como
ADR ou na seção de comparativos do \texttt{04-diagramas-arquitetura.md §\,8}.
\\[2pt] \textit{Responde: Luis Gustavo.}

\subsection*{Q8.4 — Como vocês garantem que a documentação não fica desatualizada?}
\textbf{A:} Documentos versionados no mesmo repo do código. Mudanças
arquiteturais relevantes exigem PR que atualize o ADR correspondente.
Não temos verificação automatizada de drift — é um próximo passo
(\textit{docs-as-code} com lint).
\\[2pt] \textit{Responde: Luis Gustavo.}

\newpage

\section*{Perguntas armadilha — atenção redobrada}

\begin{itemize}

  \item \textbf{\textquotedblleft Vocês usaram ChatGPT / Claude / IA?\textquotedblright}
        \\ \textbf{A:} Sim, declarado abertamente no \texttt{README.md}.
        IA foi usada para tarefas difíceis (processamento geoespacial,
        debugging de TypeScript, refatoração). Todo código foi
        revisado, testado e entendido pela equipe. Não somos uma equipe
        que copia e cola — usamos como pair-programmer.

  \item \textbf{\textquotedblleft Por que tão poucos testes (28)?\textquotedblright}
        \\ \textbf{A:} O número não é a métrica certa — o que importa é
        \textbf{o que está coberto}. Os 28 testes cobrem todos os casos
        de borda do processador (NaN, divisão por zero, máscara de água,
        polígono inválido). Mais testes redundantes não agregariam.

  \item \textbf{\textquotedblleft O sistema está em produção?\textquotedblright}
        \\ \textbf{A:} Não — é um MVP em localhost. Foi escolha
        consciente do escopo da disciplina. Hospedar exigiria autenticação
        e persistência, ambos no roadmap.

  \item \textbf{\textquotedblleft E se eu quiser usar com outra imagem que não Sentinel-2?\textquotedblright}
        \\ \textbf{A:} Hoje só Sentinel-2 \texttt{.SAFE}. O processador é
        agnóstico de fornecedor — qualquer raster com bandas RGB+NIR
        funcionaria — mas a camada de I/O é específica do formato \texttt{.SAFE}.
        Generalizar é refatoração de meia tarde.

  \item \textbf{\textquotedblleft Qual a diferença entre vocês e QGIS / ArcGIS?\textquotedblright}
        \\ \textbf{A:} EasyGis é \textbf{focado em um caso de uso}
        (agricultura de precisão por índices) com \textbf{zero curva de
        aprendizado} — abrir e clicar. QGIS é generalista e exige
        conhecimento de SIG. Não competimos com QGIS — somos uma porta de
        entrada para quem não vai aprender SIG.

  \item \textbf{\textquotedblleft O que vocês fariam diferente se começassem hoje?\textquotedblright}
        \\ \textbf{A:} Três coisas: (1) integrar Copernicus API desde o
        dia 1 (download manual virou gargalo); (2) CI desde o dia 1;
        (3) persistir talhões em SQLite — barato e adiciona valor de
        verdade pro usuário.

\end{itemize}

\vspace{8pt}
\textbf{Lembrete final:} se uma pergunta cair fora do que está aqui, a
melhor resposta é \textbf{\textquotedblleft Excelente pergunta — está
documentada em [doc X], mas posso voltar com a resposta após a
sessão.\textquotedblright}\ Nunca chute.

\end{document}
